\documentclass[letterpaper,10pt,conference]{ieeeconf}

\IEEEoverridecommandlockouts
\overrideIEEEmargins

\usepackage[T1]{fontenc}
\usepackage{amsmath}
\usepackage{booktabs}
\usepackage{graphicx}
\usepackage{newtxtext,newtxmath}
\usepackage{tabularx}
\usepackage{url}
\usepackage{xcolor}

\newcommand{\method}{CLIMB}

\title{\LARGE \bf
When Failure Is Not Difficulty: Screening Reference--Physics Misalignment\\
and Testing Adaptive Allocation on Exact Support for Humanoid Motion Tracking
}

\author{Anonymous Authors}

\begin{document}
\maketitle
\thispagestyle{empty}
\pagestyle{empty}

\begin{abstract}
Persistent tracking failure need not identify useful practice: a retargeted reference may
demand support unavailable under the declared robot and scene. We study this
reference--physics misalignment (RPM) through model-relative screening and controlled
allocation experiments. In a three-seed Unitree G1 campaign, top-1 exposure peaks at
87--89\%, an unsupported kneel-and-crawl reference repeatedly attracts practice, and
uniform sampling yields higher held-out survival in each seed. A final-trajectory
contact-capacity screen flags 2,442 of 10,705 AMASS-derived references in one retargeting
pipeline and 7 of 4,950 in a separate production pairing; two implementations agree on
39 of 40 enriched-panel decisions. Feasibility features improve cross-policy difficulty
transfer from Spearman correlation 0.567 to 0.609 over 100 held-out clips. We then bind
training to 1,184 exact temporal units and compare four allocators at equal support and
49.2 million transitions per policy. Relative-progress allocation changes exposure but
does not establish the registered tracking benefit: feasible-hard R$-$U TrackingScore
is $-0.0151$, with paired seed-level 95\% interval $[-0.0944,+0.0642]$. The broad
non-regression criterion also fails, and exploratory learning curves favor uniform.
These results separate reference admission, sampling intervention, and policy utility;
they establish neither allocation equivalence nor a general advantage of gating.
\end{abstract}

\section{Introduction}

Large retargeted motion banks have made generalist humanoid tracking a practical training recipe:
map human motion to a robot, train a policy over the resulting references in massively parallel
simulation, and spend additional updates on motions the policy currently fails. Systems such as
PHC, MaskedMimic, H2O, ExBody, BeyondMimic, and SONIC demonstrate the scale of this
recipe~\cite{luo2023phc,tessler2024maskedmimic,he2024h2o,cheng2024exbody,liao2025beyondmimic,luo2025sonic}.
Failure, completion, or recent learning progress then appears to identify where the next policy
update is most useful~\cite{liao2025beyondmimic,chen2025gmt,yang2025egm}.

That interpretation assumes every reference defines a solvable control problem. A retargeted
trajectory may instead have no admissible contact capable of supplying its demanded base wrench,
or may require joint effort beyond the modeled actuator range. We call this
\emph{reference--physics misalignment} (RPM). Under RPM, failures can reflect reference inadmissibility as well as limited
experience or control capability. Our implemented screen addresses the first distinction.
A coverage floor cannot create a behavior absent from the bank, and a progress estimate
does not identify intrinsic difficulty or forgetting. The experimental question is therefore:
given fixed reference support and a training budget, does learner-dependent allocation
produce better tracking on the same held-out objective?



We observed the consequence in a 100-motion Unitree G1 campaign. A BeyondMimic-style sampler
concentrated on the same kneel-and-crawl reference in all three seeds. Campaign peak top-1 mass was
0.884, 0.870, and 0.893; another clip owns the maximum in two seeds, but the same kneel-and-crawl
reference is the recurrent attractor. It passes ordinary kinematic checks. Yet during its
0.75--1.75\,s descent, the nearest collision geometry remains 7--10\,cm above the floor while a
median 329\,N of support demand---approximately the 327\,N model weight---has no modeled source.

\method{} makes reference admission and practice allocation separately testable
(Fig.~\ref{fig:interface}). The \texttt{refeas} screen evaluates final robot-space references
under a declared model; temporal units enforce support throughout each trial. A bounded
allocator may use policy outcomes to change exposure within that support. Whether this
feedback improves control is an empirical question, not a property of the interface.

Our contributions are:
\begin{enumerate}
  \item an RPM diagnosis and model-relative contact-capacity screen, evaluated at bank
  scale and across implementations, with feasibility features that transfer difficulty
  information between policies sharing a learner;
  \item an exact temporal-support interface and a matched four-arm, three-seed allocation
  study with a complete, inconclusive policy result and explicit exposure measurements.
\end{enumerate}
We also report a bounded repair alternative: 22/26 screened candidates qualify, yet a
one-policy paired deployment comparison does not establish an aggregate tracking gain.
Admission's learning value requires a separate on/off experiment under a common allocator;
the present allocation comparison cannot answer it.

\section{Related Work}

\subsection{Motion filtering and dynamic retargeting}

Policy-based curation already removes poorly tracked references. H2O retains AMASS motions that a
privileged imitator can follow, and ExBody2 uses an initial policy's sequence errors to select a
feasible and diverse subset~\cite{he2024h2o,ji2024exbody2}. These scores combine reference quality
with policy capability. KungfuBot instead uses a human-space center-of-mass/center-of-pressure
heuristic before robot retargeting~\cite{xie2025kungfubot}. LIMMT combines target-robot motion
heuristics, diversity, and complexity, with physical-score weights calibrated through repeated
policy training~\cite{guan2026limmt}. PHUMA filters source artifacts and retargets with joint-limit,
ground-contact, and anti-skating losses~\cite{lee2025phuma}. \method{} builds on this line but
changes the training interface: the final robot-space screen determines the support on which an
adaptive allocator is permitted to act.

Optimization-based retargeters address a complementary problem. Kinodynamic Motion Retargeting
uses rigid-body and contact constraints to construct viable locomotion, while Direct Dynamic
Retargeting uses simulator-in-the-loop optimization without an infeasible geometric
intermediate~\cite{zhang2026kinodynamic,roux2026direct}. AMO and SPIDER likewise generate viable
references through trajectory optimization or physics-based sampling with contact
guidance~\cite{li2025amo,pan2025spider}. GMR further shows that retargeting artifacts affect
downstream tracking robustness~\cite{araujo2025retargeting}. An earlier analytic study estimates
velocity, acceleration, and torque requirements for 40 upper-body motions, but does not test
whole-body environmental support~\cite{klas2023actuator}. Repair can therefore be preferable to
exclusion. Our narrower question is diagnostic: after retargeting, can the demanded robot-space
wrench be supplied by admissible contacts within modeled actuator limits?

\subsection{Adaptive allocation and evaluation}

Prioritized experience and level replay formalize the benefit and bias of nonuniform training
distributions~\cite{schaul2016per,jiang2021plr}. In humanoid tracking, BeyondMimic uses a
failure-EMA bin sampler, while GMT and EGM reweight motions or segments using tracking
outcomes~\cite{liao2025beyondmimic,chen2025gmt,yang2025egm}. Those systems combine allocation with
curation, staging, clipping, or architectural changes. Empirical failure, completion, and tracking
error identify hard examples, but do not by themselves determine whether a reference is
dynamically admissible. \method{} inserts that test before outcome-based allocation; its matched
comparison then changes only allocation over a shared set of legal trials.
Learning-progress curricula already include ALP-GMM and portable implementations
such as Syllabus~\cite{portelas2019teacher,sullivan2025syllabus}. The contribution
is the temporal-support contract and controlled test of policy utility, not
normalization of a progress score alone.

Robot-learning comparisons are sensitive to seed count, aggregation, and interval
construction~\cite{agarwal2021precipice}. Tracking adds a conditioning problem: kinematic error
among survivors can improve when a method terminates earlier. HumanTracker similarly finds that
kinematic metrics can miss support and contact failures~\cite{liu2026humantracker}. We therefore
use a liveness-weighted primary score and report common-survivor quality only beside survival.

\section{Screening and Exact Temporal Support}

\subsection{Model-relative feasibility}

Let a retargeted reference provide generalized position and velocity
$(q_t,\dot q_t)$ for the G1 at frame $t$. We apply a centered five-frame moving average to
velocity, differentiate at the reference frame rate, and smooth acceleration with the same
window. Feasibility is defined relative to a robot model, scene, collision geometry, friction,
actuator ranges, and reference. The present instantiation uses the G1 and a flat plane; electrical,
thermal, compliance, latency, and structural limits remain outside this model-relative label.

We use the compiled MuJoCo model underlying the tracking environment. Contact-free inverse
dynamics returns the generalized force required for the prescribed acceleration. The six
free-joint coordinates define the base wrench $\mathbf W_t$ that the environment must supply; the
remaining coordinates define contact-free joint torque $\boldsymbol\tau_{\mathrm{free},t}$. A collision
geometry becomes a candidate contact when its exact distance to the plane is at most
$g=0.06$\,m.

The 6\,cm band is a declared tolerance. At 3\,cm a known feasible control is flagged on 42\% of
frames because the bank carries an approximately 3\,cm ground-alignment offset. At 10\,cm the
attractor's 7--10\,cm floating descent is incorrectly granted a contact candidate. The chosen
setting lies between these observed failure modes; it is not a universal constant.

\subsection{\texttt{refeas}: contact-capacity screen}

For each candidate point we construct a four-edge pyramidal friction cone with coefficient 0.6 in
the primary model; a simulator-matched view makes non-foot contacts frictionless. Mapping
nonnegative generator magnitudes $f_t$ through the contact Jacobian yields both base-wrench
matrix $\mathbf A_t\mathbf f_t$ and joint-torque contribution $\mathbf J_t\mathbf f_t$. A weighted nonnegative least-squares
problem first exposes the unconstrained residual, with angular-wrench rows weighted by two. We
then solve a linear program minimizing $\ell_1$ wrench slack subject to
\begin{equation}
-\boldsymbol\tau_{\max} \leq \boldsymbol\tau_{\mathrm{free},t}-\mathbf J_t\mathbf f_t \leq \boldsymbol\tau_{\max},
\qquad \mathbf f_t\geq0.
\label{eq:contact}
\end{equation}
The translational component of the resulting weighted residual defines unsupported force. With
no candidate contact, the residual is the contact-free demand.

We report the fraction of frames with no contact candidate (\texttt{airborne\_frac}), the fraction
whose torque-limited unsupported force exceeds half the robot weight
(\texttt{infeasible\_frac}), normalized unsupported-force integral, and joint-torque
infeasibility. A clip crosses the primary rule only when
\texttt{infeasible\_frac}$>0.10$ (strict inequality). Ballistic flight is not flagged merely for
being airborne: near-gravity acceleration has little external support demand. Conversely,
nonballistic hovering or descent without support is flagged.

The screen takes approximately one CPU-second per clip in the primary implementation. An
independent implementation processes the production bank at 0.145 CPU-s per clip
(0.84\,ms/frame). The screen therefore runs as an offline bank-ingestion stage rather than in the
policy control loop.

\subsection{Repair as a routing alternative}
DFRP applies a bounded contact projection to qualified floating-reference cases. It lowers
and smooths root clearance, then adjusts supporting-leg joints with damped contact inverse
kinematics and joint-limit projection. The changed trajectory is re-screened and must
satisfy residual infeasibility at most 5\%, root displacement at most 8\,cm, valid joint
limits, contact-IK residual at most 10\,mm, and at least one legal 50-step start. Missing-scene
cases bypass projection. This is an implemented routing option, not a guarantee of usable
control: repair changes the tracking target and needs its own paired policy evaluation.

\subsection{Exact-support gated allocation}

Clip classification is too coarse because one reference can contain both supportable and
unsupported intervals. We reduce per-frame screens to ordered feasible runs, retain only 50-step
windows wholly inside a run, and bind each motion and sidecar by SHA-256. Across 800 training
clips, 1,841 source runs yield 1,184 admissible units and 368,951 legal starts after 657 short runs
are discarded. A unit is an attribution stratum; the deployment prior is uniform over legal
starts, not units or clips.

For an admitted interval $\mathcal A_u$, define legal starts by every reference
frame accessed during horizon $H$ and lookahead configuration $\mathcal L$:
\begin{equation}
\mathcal S_u=\{s:s+\mathcal K(H,\mathcal L)\subseteq\mathcal A_u\}.
\end{equation}
The current interface includes the terminal observation and uses current-frame reference
observations. ``Exact'' denotes membership relative to this screen and access contract;
it is not a certificate of physical feasibility.

Let $b_u$ be the normalized legal-start prior over admitted units and
$g_u=|s_u(k)-s_u(k-10)|$ the nonnegative change of Beta-smoothed conditional success
at sampler clock $k$. Uniform U uses $b$. Absolute-progress A retains the earlier
fixed additive floor. Relative-progress R, before active caps and for complete
history with $\mu=\sum_u b_ug_u>0$, is exactly
\begin{equation}
p_u=0.8b_u+0.2\frac{b_ug_u}{\mu}.
\label{eq:gate}
\end{equation}
The implemented prior/cap fallback covers incomplete history or zero $\mu$.
Common positive rescaling of $g$ leaves Eq.~\eqref{eq:gate} unchanged, and its
pre-cap total variation from $b$ is at most 0.2. These are exposure properties,
not guarantees of useful practice. D uses conditional failure under the same
support; its startup and floor/cap composition differ from R, so R$-$D compares
allocator designs rather than isolating ranking alone.

Intervals do not all receive the same remedy. Admissible runs and DFRP-qualified repairs enter the
exact support; missing-context and residual failures do not. This routing preserves feasible
subsequences rather than treating every screen flag as a whole-clip deletion instruction.
The four-arm comparison uses the original admitted bank; it substitutes no repaired targets.

The runtime samples only legal starts, emits explicit truncation at a segment boundary, never
wraps into a rejected frame, and serializes its deterministic generator and sufficient
statistics. Per-unit and per-clip caps bound concentration. Missing sidecars, changed motion
hashes, invalid starts, invalid frames, or censored resets fail the contract.

The paired evaluator freezes the motion, start, replicate, environment-noise seed, and
dynamics-noise seed across policies. It assigns the reference before the first observation,
disables auto-reset until terminal channels are captured, rejects invalid offsets, and records
checkpoint, task, condition, reference, evaluator, and output hashes. The 100-clip panel is
name- and hash-disjoint from training and contains 2,800 paired conditions.

\begin{figure*}[t]
  \centering
  \includegraphics[width=\textwidth]{f1_evidence_interface.pdf}
  \caption{\textbf{Admission, allocation and policy utility are separate tests.}
  The screen supplies model-relative support; the temporal interface enforces legal
  reference access; independent evaluation measures the resulting policy. The tested
  outcome-to-allocation feedback changed exposure but did not establish the registered
  tracking benefit. Neither screen qualification nor exact support implies hardware
  feasibility. Repair is a separate, target-changing route.}
  \label{fig:interface}
\end{figure*}

\section{Experimental Design}

\subsection{E1: RPM attractor diagnosis}

We analyze the previously sealed 100-motion campaign: 4,096 environments, 4,000 PPO iterations,
three seeds per arm, and 100 disjoint evaluation motions with eight episodes per clip. It includes
failure-adaptive, clip-uniform, and normalized grounded samplers. We report exposure, attractor
identity, held-out survival, and the attractor's modeled contact-capacity residual. This campaign
establishes the motivating failure case; the later four-arm study separately compares exact-support allocators.
The historical adaptive arm adds exploration mass before normalization, which does
not guarantee a probability floor. Its collapse therefore diagnoses that particular
sampler/reference combination. The current D baseline instead retains a normalized
80\% prior component; the historical result does not establish that every
failure-adaptive allocator will collapse under RPM.


\subsection{E2: bank scale, agreement, and transfer}

The primary corpus contains 10,705 AMASS-derived references retargeted through
\texttt{whole\_body\_tracking} to G1. We report raw counts under the fixed rule, category/source
breakdowns, and contamination in the historical evaluator. A separate experiment covers all
4,950 references in the filtered BONES-SEED production pairing. We never pool their prevalence.
A deterministic 40-clip panel, 20 from each bank and enriched for flags, passes through both
implementations to measure score and decision agreement; it is not a prevalence sample.
Separately, a ridge model fits 11 intrinsic reference features to per-clip difficulty under one
curriculum and ranks difficulty under another on the same 100 held-out clips. We compare intrinsic
features with and without the three \texttt{refeas} features against 200 random three-feature
additions; this is cross-policy transfer, not cross-architecture transfer.

\subsection{E3: DFRP repair qualification}

We freeze a source- and severity-stratified CPU panel before repair: 26 strict-flagged candidates
span three root-displacement bands and two initial-infeasibility bands, and four source-matched
feasible clips serve as no-op controls. The primary outcome is exact-ready flagged count under the
5\% residual, 8\,cm root, joint-limit, 10\,mm contact-IK, source-hash, and legal-start gates. The
gate requires at least 20/26 candidates, 4/4 byte-identical ready controls, zero integrity or
qualification failures among admitted clips, and a reproducible manifest payload. Runtime and
root, joint, body, velocity, and acceleration deviations are secondary diagnostics.

\subsection{Four-arm allocation on unchanged exact support}
An earlier two-arm E4 protocol selected absolute-progress settings using allocation-only
calibration. Its long-run manipulation check subsequently failed, leaving policy endpoints
closed. A separate prospective confirmation retained A as a declared comparator and added
relative progress R and conditional failure D alongside uniform U.

The completed study uses seeds 21, 22 and 23, 512 environments, 4,000 PPO iterations and
24 rollout steps: 49,152,000 simulator transitions per policy. Support, learner, rewards
and reference targets are common. All twelve training gates passed before any held-out
cell opened. Checkpoints 1000, 2000, 3000 and 3999 give 48 evaluation cells on 100 unchanged
held-out clips, including 25 feasible-hard clips selected from reference features. Each cell
contains 2,800 conditions, evaluated over uninterrupted windows of up to three seconds.

The actor and critic are observation-normalized ELU multilayer perceptrons with
hidden widths 512/256/128, taking 160 and 286 observation coordinates, respectively;
the actor outputs 29 joint-position commands. Physics runs at 200\,Hz with 50\,Hz
control. PPO uses Adam (initial learning rate $10^{-3}$, adaptive KL target 0.01),
five epochs, four minibatches, clipping 0.2, discount 0.99, GAE coefficient 0.95
and entropy coefficient 0.005. These settings are shared across arms.

The primary score is liveness-weighted reference tracking:
\begin{equation}
h\exp\!\left(-\frac{\mathrm{MPKPE}}{0.30\,\mathrm{m}}
             -\frac{\mathrm{anchor\ error}}{0.40\,\mathrm{rad}}\right),
\end{equation}
where $h$ is survived-horizon fraction, MPKPE is root-relative body-position error,
and anchor error is orientation error. We average conditions within clip, clips within
panel and paired differences across training seeds. Independent replication is the
trained seed pair, not the number of conditions.

The frozen primary decision requires final hard R$-$U mean at least $+0.02$, a positive
lower two-sided 95\% paired seed t bound (df=2), and all-panel lower bound above $-0.01$.
A pass would establish a positive effect with a target-reaching point estimate, not a
true effect of at least $+0.02$. Paired hierarchical bootstrap intervals are supplementary.
Learning curves, normalized area under the observed curve (AULC), and attainment of the
paired U final score are exploratory; non-attainment is retained. No seeds are appended
to the completed decision.

\section{Results}

\subsection{E1: Curriculum collapse under RPM}

Adaptive peak top-1 mass is 0.884/0.870/0.893 across campaign seeds, and the same BMLmovi
kneel-and-crawl clip repeatedly becomes dominant in all three. The maximum belongs to another
clip in two seeds. Uniform sampling achieves higher held-out endpoint survival in every paired
seed: differences (uniform minus adaptive) are $+0.0300$, $+0.0275$, and $+0.0300$. With only
three seeds, the sign pattern is the useful description; the minimum attainable paired
permutation $p$-value is 0.125.

The reference has no joint-limit violation and peak joint speed is at most 5.6\,rad/s. The
dynamic screen instead localizes the failure. From 0.75--1.75\,s, 86\%
of frames have no candidate contact and median unsupported force is 329\,N; the model weighs
327\,N. The 1.75--7.25\,s kneeling interval is supportable through non-foot contacts, and the
8.0--8.5\,s rise again becomes unsupported. Every trained policy has zero survival from frame
zero, while the tested late 8\,s offset survives.

\subsection{E2: Bank-scale screening and difficulty transfer}

Under the strict rule, 2,442 of 10,705 primary-pipeline clips are flagged (22.8\%). Category
rates range from 12.9\% for quiet motion to 58.6\% for dynamic motion, and source rates from
0.1\% for GRAB to 100\% for CNRS. Severity-stratified inspection shows why labels are
insufficient: five inspected CNRS cases are ordinary fast walks with extended floating
intervals, while the Transitions sample mixes ingest defects with acrobatic content.

Only 7 of 4,950 production clips (0.14\%) cross the same nominal rule, while 111 (2.24\%) are
airborne on more than 10\% of frames. Five of seven flags are jumps, four naming a 50\,cm box
absent from the declared flat scene. Their verdict is \emph{unsupported on this plane}, routing them
toward missing terrain or exclusion rather than geometric root projection.

Across the stratified 40-clip panel, \texttt{infeasible\_frac} rank correlation is 0.9836,
\texttt{airborne\_frac} correlation is 0.9974, and strict decisions agree for 39/40 clips
(Cohen's $\kappa=0.9485$). The disagreement is \texttt{burpee\_002\_\_A362\_M}
(0.019 versus 0.136). Figure~\ref{fig:bank} keeps the full-bank denominators separate and uses
only this same-input panel to check implementation agreement; it does not establish corpus or
retargeter equivalence.

The difficulty decomposition is also visible across policies trained with different curriculum
allocations. On the same 100 held-out clips, an 11-feature intrinsic atlas fit to adaptive-policy
difficulty ranks uniform-policy difficulty at $\rho=0.567$. Adding
\texttt{infeasible\_frac}, \texttt{airborne\_frac}, and normalized unsupported impulse raises
the rank correlation to $0.609$, above 198 of 200 random three-feature additions (one-sided
$p=0.010$). All six directed policy pairs move in the same direction and four have $p<0.05$;
the two pairs targeting the grounded policy do not. This is transfer across policies sharing an
architecture, not evidence of cross-architecture or cross-robot transfer.

\begin{figure*}[t]
  \centering
  \includegraphics[width=\textwidth]{f2_bank_scale.pdf}
  \caption{Bank-scale screen and implementation agreement. (a) Bars are normalized within two
  separate corpus/pipeline denominators and apply the strict
  \texttt{infeasible\_frac}$>0.10$ rule: 2,442/10,705 for AMASS$\rightarrow$WBT$\rightarrow$G1
  and 7/4,950 for filtered BONES-SEED$\rightarrow$G1. Corpus, filtering, robot file, friction,
  and implementation differ, so this is not a causal retargeter comparison. (b) On a
  flag-enriched, deterministic 20+20 same-clip panel, strict decisions agree for 39/40 clips
  (Spearman $\rho=0.9836$, Cohen's $\kappa=0.9485$). That panel checks implementation agreement;
  its stratified selection is not a prevalence sample.}
  \label{fig:bank}
\end{figure*}

\subsection{E3: Contact-manifold repair via DFRP}

DFRP qualifies 22/26 flagged candidates (84.6\%) and all four feasible controls pass as
byte-identical no-ops. The resulting 26-clip curated view contains 36 admissible units and 10,561
legal 50-step starts. Median and 95th-percentile CPU runtime are 2.57 and 6.29\,s per clip. Two
candidates remain above the 5\% residual-infeasibility ceiling, and two more reach that ceiling
but violate the 10\,mm contact-IK bound; all four are quarantined. Thus 84.6\% is a qualification
rate on the frozen panel, not a bank-wide recovery rate or a policy-benefit estimate.

\begin{table}[t]
\centering
\caption{Repair qualification and deployment answer different questions.}
\label{tab:repair}
\begin{tabularx}{\columnwidth}{lX}
\toprule
Evidence / status & Result and scope \\
\midrule
Qualification / measured & 22/26 flagged candidates; 4/4 feasible controls byte-identical. \\
Deployment / exploratory & One policy, 26 clips, 656 paired conditions per arm: raw 0.3925,
repaired 0.3915; difference $-0.0010$, clip-bootstrap 95\% CI $[-0.0086,+0.0080]$. \\
\bottomrule
\end{tabularx}
\end{table}
The fixed-policy deployment result (Table~\ref{tab:repair}) does not establish an
aggregate gain, safety, or equivalence. It compares tracking of raw and repaired targets;
it is not a retraining effect or improved tracking of the unchanged original reference.

\subsection{Completed allocation result: inconclusive}
\label{sec:e4results}

The earlier E4 absolute-progress arm reaches mean post-warm-up TV 0.0297, below its
0.05 gate. Its disposition is \texttt{not\_tested}; it supplies no policy null.
The newer confirmation retains A under its own comparator rule. R changes exposure
in all three seeds (mean TV 0.08348/0.08249/0.08233), as does D
(0.08391/0.08587/0.08398). All twelve training runs pass their declared gates with
zero recorded invalid or censored events. Across 492 saved states, exact sampler
and checkpoint replay establishes intervention integrity.

\begin{table}[t]
\centering
\caption{Final TrackingScore differences. Three paired trained seeds;
two-sided seed t intervals, df=2.}
\label{tab:confirmation}
\begin{tabular}{lrr}
\toprule
R$-$U & Feasible-hard & All-panel \\
\midrule
Seed 21 & $-0.03400$ & $-0.04892$ \\
Seed 22 & $-0.03299$ & $-0.01402$ \\
Seed 23 & $+0.02178$ & $+0.03500$ \\
Mean & $-0.01507$ & $-0.00932$ \\
95\% lower & $-0.09436$ & $-0.11405$ \\
95\% upper & $+0.06422$ & $+0.09542$ \\
\bottomrule
\end{tabular}
\end{table}

The frozen primary result is \emph{inconclusive} (Table~\ref{tab:confirmation}). Neither
the primary rule nor the broad non-regression guard passes. Two hard-panel seeds favor U
and one favors R. The interval spans both meaningful harm and benefit; it does not establish
either, nor does failure of the guard prove regression. The supplementary paired bootstrap
hard interval is $[-0.04417,+0.02532]$ and cannot override this decision.
Descriptive final hard R$-$D is $-0.00675$ with interval $[-0.02189,+0.00840]$;
R$-$A is $-0.00668$ with interval $[-0.07243,+0.05907]$. These contrasts do not
establish equivalence or progress-specific superiority.

\begin{figure*}[t]
\centering
\includegraphics[width=\textwidth]{paired_results_and_learning_curves.pdf}
\caption{Completed allocation confirmation. Seed pairs, means and two-sided 95\% t
intervals appear at left and center. Dashed lines are the hard point-estimate target and
all-panel lower-bound margin. Thin learning curves are seeds, thick curves are means;
the cost axis is simulator transitions. Curves and attainment are exploratory.}
\label{fig:confirmation}
\end{figure*}

Exploratory normalized hard AULC R$-$U is negative in all three seeds:
$-0.03978$, $-0.03042$, $-0.01676$ (mean $-0.02898$; descriptive 95\% t interval
$[-0.05774,-0.00023]$). Two R seeds never attain their paired U final score on the
observed grid. These data do not support an efficiency-accelerator claim; sparse checkpoints
do not identify a precise crossover. The secondary interval is not a new confirmatory
harm test. Figure~\ref{fig:confirmation} reports the complete curves.

\paragraph{Earlier controls remain separate}
Table~\ref{tab:controls} separates prior protocol statuses. None is pooled with the
completed four-arm result or substituted for a matched admission on/off test.
\begin{table}[t]
\centering
\caption{Earlier controls, with their original scope and status.}
\label{tab:controls}
\begin{tabularx}{\columnwidth}{lX}
\toprule
Study / status & Bounded finding \\
\midrule
E-HYG / sealed & Whole-clip pruning: feasible-panel survival $0.918\to0.907$;
$\Delta=-0.0101$, one-sided permutation $p=0.951$. \\
Soft weighting / failed gate & Hard-rejected mass 0.199; this does not implement exact admission. \\
E4 / not tested & Mean TV 0.0297 below 0.05; policy endpoint unopened. \\
Repair-all / exploratory & Deployment contrast $+0.0397$ misses $+0.05$ target and coverage gate. \\
\bottomrule
\end{tabularx}
\end{table}
Segmentation retained 12.5 of the 20.2 minutes in 99 flagged training clips and
lost only three whole clips at zero guard. Retained duration is not demonstrated
retained skill. A matched on/off admission comparison under D is the next distinct test;
no incomplete result enters this manuscript.

\section{Limitations}

\method{}'s admission label is model-relative. The current screen instantiates one G1 model and a
flat scene with collision contacts, friction, and actuator force limits. Hardware deployment would
add electrical, thermal, compliance, latency, and structural constraints and then require
closed-loop validation. Terrain-bearing references likewise require their terrain in the scene;
the four box-jump flags in the production bank are routed to missing context rather than declared
bad motion.

Finite differences, the 6\,cm contact band, and the half-weight residual threshold introduce
modeling choices. Local sensitivity checks bound the two thresholds on this bank, but a new robot
or retargeter must recalibrate them. The 22.8\% and 0.14\% rates therefore remain separate
corpus--pipeline measurements. A smooth feasibility weight derived from residual slack could
retain borderline cases, but would couple admissibility and allocation; it requires its own
matched ablation against the current binary gate.

The allocation study has only three independent seed pairs. Its observed hard-panel
standard deviation 0.03192 yields a 95\% t half-width 0.07929, much larger than the
$+0.02$ point target. This is an observed precision limitation, not a retrospective
power guarantee or evidence of equivalence. Thousands of evaluation conditions cannot
replace independent training. The current result neither establishes benefit nor rules
out a target-sized benefit; a future replication requires its own fixed sample size.

The current simulator has knee effort clamps of approximately $\pm139$\,N\,m,
zero added command delay and plane terrain; physical G1 transfer is untested. Numerical
parity remains unresolved in a separate physical-sensitivity development check, so
its perturbation scores are not reported as evidence. A continuous adapter traverses
8.58-second training references, but comparative full-motion execution is untested.
The frozen-policy instrumentation pilot does not establish a noise floor; its planned
17.6-million-transition burn-in makes it a separate study. Practice interventions,
reliability-calibrated progress and hardware evaluation remain future work, not
explanations established by the current allocation result.

\section{Conclusion}
RPM makes persistent failure an ambiguous instruction for practice. The diagnosed
campaign, bank-scale screen and cross-policy difficulty transfer motivate explicit
measurement of reference admissibility. Exact temporal support then permits a controlled
allocation test: relative progress changes exposure, but the completed three-seed
comparison does not establish its registered tracking benefit. A qualified repair likewise
does not establish a fixed-policy gain. Reference qualification, exposure and control
utility therefore need separate evidence. The learning value of admission under a common
allocator remains the next direct experiment.

\section*{Reproducibility and Acknowledgments}
The stack uses mjlab v1.6.0 at commit \texttt{0fb8a681136b} and MuJoCo 3.11.0.
Motion, sidecar, unit-table, task, checkpoint, evaluator, condition and output identities
are SHA-256 bound. Licensed motion payloads cannot be redistributed; release requires
reconstruction instructions and permitted aggregate metadata. OpenAI Codex and
Anthropic Claude assisted with code, analysis tooling, figures and manuscript drafting
in Sections I--VII; claims are tied to verified artifacts, with final author review
required before submission.

\bibliographystyle{IEEEtran}
\bibliography{references}

\end{document}
